\section{日々の変動と、シミュレーションの当たりやすさ}
\label{sec:fidelity}

ここまでの章は「いつ・どこで・どの高度の風が、ふつうどちらへどれだけ吹くか」を
まとめてきた。本章はそこから一歩進んで、打ち上げを計画する人が実際に抱える
2 つの問いに答える。

\begin{quote}
\begin{enumerate}
\item \textbf{日ごとの風向・風速の変わり方は季節によって違うのか}。
\item \textbf{シミュレーション通りになりやすいのは、どの季節か}。
\end{enumerate}
\end{quote}

3 領域とも同じ手順で測った。
北海道と南の 2 領域（和歌山・愛媛）では季節の効き方がまるで違うので、
片方だけを見て一般論を言うことができない。
一方、和歌山と愛媛はほとんど同じ結果になった。

1 つ目はデータを見れば答えが出る。一方で 2 つ目は、
「シミュレーション通り」が何を指すのかを決めないと答えようがない。
そこで先に、落下点の予測が外れる原因を整理しておく。

\subsection{予測が外れる原因は 4 つあり、測れるのは 3 つ}
\label{sec:fidelity-scope}

\begin{table}[H]
\centering
\caption{落下点の予測を外す 4 つの原因と、再解析から測れるかどうか}
\label{tab:fidelity-sources}
\begin{tabular}{@{}llc@{}}
\toprule
 & 原因 & 再解析から測れるか \\
\midrule
(1) & シミュレータが風を\textbf{打ち上げ時刻}で固定している & 測れる \\
(2) & シミュレータが風を\textbf{打ち上げ地点}で固定している & 測れる \\
(3) & 打ち上げが予定から N 時間ずれる & 測れる \\
(4) & 与えた\textbf{予報自体}が外れている & 測れない \\
\bottomrule
\end{tabular}
\end{table}

(4) を測るには「過去に出された予報」の保管庫が要る。
本研究が持っているのは「過去に実際どうだったか」の再解析であって、
予報の保管庫ではない。したがって (4) は本章の対象外である。

この点は結論の読み方に直結するので、先に言っておく。
実務では (4) がいちばん大きい誤差である可能性が高い。
本章が出す数字は、\textbf{予報が完璧だったとしてもこれだけは外れる}という下限であって、
実際の外れ幅の見積もりではない。

\subsection{シミュレータが置いている 2 つの仮定}

\texttt{backend/bfps/core/wind.py} の \texttt{WindProfile.wind\_at} は、
緯度・経度・時刻を引数で受け取りながら、そのすべてを無視して、
打ち上げ点で 1 回取得したプロファイルを返す。
呼び出し側はそれを飛行全体（約 \SI{137}{\minute}）に使い回す。

その間に機体は水平に \SIrange{60}{300}{\kilo\metre} 移動し、風そのものも変わる。
つまり機体は、もう自分がいない場所の、もう過ぎた時刻の風で飛ばされている。

\begin{whybox}{これは設計ではなく、風の取り出し方だけの問題である}
\texttt{trajectory.py} の \texttt{horizontal\_derivative} は、
時刻 \texttt{tt} を引数で受け取っている。
RK4 の呼び出し側も各段で \texttt{t}、\texttt{t+dt/2}、\texttt{t+dt} をきちんと渡している。
ところが関数の中身は \texttt{tt} を捨てて、打ち上げ時刻を使う。

位置のほうも同じで、RK4 は緯度経度を正しく更新して渡しているのに、
受け取った \texttt{WindProfile} がそれを見ない。

積分の骨組みは 4 次元の風を扱えるようにできていて、
風を返す側だけがそうなっていない。
直すこと自体は難しくない。難しいのは\textbf{直す価値があるかどうかの判断}で、
そのためには誤差の大きさを知る必要がある。本章はその数字を出す。
\end{whybox}

\subsection{日々の変化は季節でどう違うか}
\label{sec:fidelity-variability}

\subsubsection{風はどれだけ動くか}

まず、時間 $L$ が経ったときに風がどれだけ変わるかを見る。
測るのはベクトルの差 $|\mathbf{V}(t+L) - \mathbf{V}(t)|$ である。

\begin{whybox}{なぜ風速の差ではなくベクトルの差なのか}
風速だけを比べると、強さを保ったまま \ang{90} 向きを変えた風が
「まったく変わっていない」と判定されてしまう。
その風の中で上げたスペースバルーンは、まるきり違うところに落ちる。

向きの変化と強さの変化を一度に扱うには、ベクトルの差をとるしかない。
なお、角度そのものを引き算してはいけない理由は
\ref{sec:circular} 節に述べたとおりで、ここでも角度は一度も引いていない。
\end{whybox}

季節ごとの組は、両端が同じ季節に入るものだけを数えている。
遅れ 10 日では 2 月末に始まった組が 3 月に着いてしまうので、
そのまま数えると、季節差がいちばん出てほしい長い遅れのところで
季節どうしが混ざってしまう。

\begin{figure}[H]
\centering
\includegraphics[width=1.0\linewidth]{fidelity/G1-wakayama_change_curve.pdf}
\caption[風の変化量]{遅れ $L$ を置いたときの風の変化量 $|\mathbf{V}(t+L)-\mathbf{V}(t)|$。実線が中央値、帯が p90 まで。横軸は対数で、6 時間から 20 日まで。組は両端が同じ季節に入るものだけを数えている。\SI{10}{\hecto\pascal} と \SI{30}{\hecto\pascal} の曲線が遅れ 24 時間でいったん下がるのは大気潮汐による見かけの効果である。}
\label{fig:g1-wakayama}
\end{figure}

\begin{figure}[H]
\centering
\includegraphics[width=1.0\linewidth]{fidelity/G1-hokkaido_change_curve.pdf}
\caption[風の変化量]{遅れ $L$ を置いたときの風の変化量 $|\mathbf{V}(t+L)-\mathbf{V}(t)|$。実線が中央値、帯が p90 まで。横軸は対数で、6 時間から 20 日まで。組は両端が同じ季節に入るものだけを数えている。\SI{10}{\hecto\pascal} と \SI{30}{\hecto\pascal} の曲線が遅れ 24 時間でいったん下がるのは大気潮汐による見かけの効果である。}
\label{fig:g1-hokkaido}
\end{figure}

\begin{figure}[H]
\centering
\includegraphics[width=1.0\linewidth]{fidelity/G1-ehime_change_curve.pdf}
\caption[風の変化量]{遅れ $L$ を置いたときの風の変化量 $|\mathbf{V}(t+L)-\mathbf{V}(t)|$。実線が中央値、帯が p90 まで。横軸は対数で、6 時間から 20 日まで。組は両端が同じ季節に入るものだけを数えている。\SI{10}{\hecto\pascal} と \SI{30}{\hecto\pascal} の曲線が遅れ 24 時間でいったん下がるのは大気潮汐による見かけの効果である。}
\label{fig:g1-ehime}
\end{figure}


和歌山の \SI{200}{\hecto\pascal} では、遅れ 1 日で春 \SI{16.8}{\metre\per\second}、
夏 \SI{11.7}{\metre\per\second} と開く。
\textbf{春がもっとも荒く、夏がもっとも穏やか}である。
北海道でも季節の順序は同じだが、季節間の開きが小さい。
愛媛の曲線は和歌山とほとんど重なる。

\SI{10}{\hecto\pascal} と \SI{30}{\hecto\pascal} の曲線は、
遅れ 12 時間から 24 時間にかけていったん下がる。
これは大気潮汐で、1 日周期の成分が 24 時間後にはほぼ元に戻るためである。
遅れを 1 日の整数倍にとると変化が小さく見える、という見かけの効果なので、
この 2 つの高度では 6 時間と 12 時間の値を読むときに注意がいる。

\subsubsection{風の偏差が続く時間}

次に、風の偏差がどれだけ持続するかを見る。
各季節・各高度で、偏差の自己相関が $1/e$ を切る遅れを求めた。
偏差はその季節自身の平均からとっている。
そうしないと、強い西風が定常に吹いているだけの季節が、
「持続性が高い」と高く出てしまう。

\begin{figure}[H]
\centering
\includegraphics[width=0.95\linewidth]{fidelity/G2-wakayama_decorrelation.pdf}
\caption[風の偏差が続く時間]{風の偏差の自己相関が $1/e$ を切るまでの日数。偏差はその季節自身の平均からとっているので、定常な風が強いだけの季節が高く出ることはない。棒がなく「20 日超」とある組は、20 日の遅れでもまだ相関が $1/e$ を上回っていた場合で、偏差が続く時間が短いのではなく長いことを意味する。}
\label{fig:g2-wakayama}
\end{figure}

\begin{figure}[H]
\centering
\includegraphics[width=0.95\linewidth]{fidelity/G2-hokkaido_decorrelation.pdf}
\caption[風の偏差が続く時間]{風の偏差の自己相関が $1/e$ を切るまでの日数。偏差はその季節自身の平均からとっているので、定常な風が強いだけの季節が高く出ることはない。棒がなく「20 日超」とある組は、20 日の遅れでもまだ相関が $1/e$ を上回っていた場合で、偏差が続く時間が短いのではなく長いことを意味する。}
\label{fig:g2-hokkaido}
\end{figure}

\begin{figure}[H]
\centering
\includegraphics[width=0.95\linewidth]{fidelity/G2-ehime_decorrelation.pdf}
\caption[風の偏差が続く時間]{風の偏差の自己相関が $1/e$ を切るまでの日数。偏差はその季節自身の平均からとっているので、定常な風が強いだけの季節が高く出ることはない。棒がなく「20 日超」とある組は、20 日の遅れでもまだ相関が $1/e$ を上回っていた場合で、偏差が続く時間が短いのではなく長いことを意味する。}
\label{fig:g2-ehime}
\end{figure}


\textbf{ここで北海道と南の 2 領域が大きく分かれる。}

\begin{table}[H]
\centering
\caption{風の偏差が $1/e$ を切るまでの日数。和歌山と愛媛は季節で 8 倍動くのに、
北海道はほとんど動かない。}
\label{tab:decorrelation-both}
\small
\setlength{\tabcolsep}{3pt}
\sisetup{table-format=2.1}
\begin{tabular}{@{}lSSSSSSSSSSSS@{}}
\toprule
& \multicolumn{4}{c}{北海道} & \multicolumn{4}{c}{和歌山} & \multicolumn{4}{c}{愛媛} \\
\cmidrule(lr){2-5}\cmidrule(lr){6-9}\cmidrule(lr){10-13}
気圧面 & {冬} & {春} & {夏} & {秋} & {冬} & {春} & {夏} & {秋} & {冬} & {春} & {夏} & {秋} \\
\midrule
\SI{10}{\hecto\pascal}  & 4.9 & 7.6 & 14.0 & 8.3  & 4.2 & 6.5 & 17.0 & 8.6  & 4.2 & 6.4 & 17.2 & 8.9 \\
\SI{30}{\hecto\pascal}  & 4.7 & 8.4 & 10.5 & 8.7  & 4.7 & 9.2 & 18.5 & 14.9 & 4.6 & 9.4 & 19.0 & 19.3 \\
\SI{100}{\hecto\pascal} & 3.3 & 2.5 & 3.5  & 4.1  & 2.4 & 3.8 & 15.4 & {20 超} & 2.3 & 4.1 & 15.5 & {20 超} \\
\SI{200}{\hecto\pascal} & 1.8 & 1.2 & 1.8  & 1.7  & 2.0 & 1.7 & 10.7 & 15.3 & 1.9 & 1.9 & 10.8 & 16.3 \\
\SI{500}{\hecto\pascal} & 1.0 & 0.9 & 1.8  & 1.1  & 1.0 & 1.1 & 4.9  & 2.5  & 1.0 & 1.1 & 5.6  & 2.8 \\
\SI{850}{\hecto\pascal} & 0.8 & 0.8 & 1.2  & 0.8  & 0.7 & 0.7 & 1.8  & 0.9  & 0.7 & 0.8 & 2.0  & 0.9 \\
\bottomrule
\end{tabular}
\end{table}

和歌山の \SI{200}{\hecto\pascal} は冬 \num{2.0} 日に対し秋 \num{15.3} 日で、
\SI{100}{\hecto\pascal} の秋にいたっては 20 日を見ても $1/e$ を切らない。
つまり\textbf{夏と秋の風はいったん決まると長く続き、冬と春の風は 1 日で入れ替わる}。
愛媛も冬 \num{1.9} 日、秋 \num{16.3} 日で、和歌山と同じ形である。

ところが北海道の \SI{200}{\hecto\pascal} は
冬 \num{1.8}、春 \num{1.2}、夏 \num{1.8}、秋 \num{1.7} 日と、
\textbf{どの季節でも 2 日を切る}。
対流圏上部については、北海道の風は季節を問わず 1--2 日で入れ替わる。

差が出るのは成層圏だけで、\SI{10}{\hecto\pascal} と \SI{30}{\hecto\pascal} の夏は
北海道でも \numrange{10}{14} 日持続する。
ただし和歌山・愛媛の \numrange{17}{19} 日よりは短い。

\begin{whybox}{なぜ北海道だけ季節差が消えるのか}
和歌山と愛媛は亜熱帯ジェットの南縁にあたり、夏にはジェットが北へ抜けて
穏やかで持続的な流れだけが残る。

北海道は寒帯前線ジェットの通り道で、夏でも低気圧が次々に通過する。
\ref{sec:hokkaido} 章で \SI{200}{\hecto\pascal} の平均風速が
夏でも \SI{29.9}{\metre\per\second} あったのはそのためで、
和歌山の夏（\SI{22.7}{\metre\per\second}）より強い。
流れが止まらないので、偏差も長く続かない。

この違いは第 \ref{sec:validation} 章の結果とも整合する。
北海道で気候値の効きが弱い（誤差削減 17 \%、和歌山は 45 \%、愛媛は 47 \%）のは、
季節が決める分が小さく日々の天気が決める分が大きいからで、
偏差が続く時間の短さも、同じ事情から来ている。
\end{whybox}

\subsubsection{ところが風向は逆の順序になる}

ここまでだけを読むと、夏は予測しやすい季節に見える。
しかし日平均風向の前日差をとると、違う結果になる。

\begin{figure}[H]
\centering
\includegraphics[width=0.95\linewidth]{fidelity/G3-wakayama_direction_change.pdf}
\caption[風向の日々の振れ]{日平均風向の前日差（符号なし）。成分に直してから差をとり、$[-180, 180]$ 度に畳んでいるので、\ang{350} と \ang{10} の差は \ang{20} になる。図\ref{fig:g2-wakayama} と季節の順序が高度によって逆転する点に注意。}
\label{fig:g3-wakayama}
\end{figure}

\begin{figure}[H]
\centering
\includegraphics[width=0.95\linewidth]{fidelity/G3-hokkaido_direction_change.pdf}
\caption[風向の日々の振れ]{日平均風向の前日差（符号なし）。成分に直してから差をとり、$[-180, 180]$ 度に畳んでいるので、\ang{350} と \ang{10} の差は \ang{20} になる。図\ref{fig:g2-hokkaido} と季節の順序が高度によって逆転する点に注意。}
\label{fig:g3-hokkaido}
\end{figure}

\begin{figure}[H]
\centering
\includegraphics[width=0.95\linewidth]{fidelity/G3-ehime_direction_change.pdf}
\caption[風向の日々の振れ]{日平均風向の前日差（符号なし）。成分に直してから差をとり、$[-180, 180]$ 度に畳んでいるので、\ang{350} と \ang{10} の差は \ang{20} になる。図\ref{fig:g2-ehime} と季節の順序が高度によって逆転する点に注意。}
\label{fig:g3-ehime}
\end{figure}


\begin{table}[H]
\centering
\caption{日平均風向の前日差（中央値、度）。高度で季節の順序が入れ替わる。}
\label{tab:direction-change}
\small
\setlength{\tabcolsep}{3pt}
\sisetup{table-format=2.1}
\begin{tabular}{@{}lSSSSSSSSSSSS@{}}
\toprule
& \multicolumn{4}{c}{北海道} & \multicolumn{4}{c}{和歌山} & \multicolumn{4}{c}{愛媛} \\
\cmidrule(lr){2-5}\cmidrule(lr){6-9}\cmidrule(lr){10-13}
気圧面 & {冬} & {春} & {夏} & {秋} & {冬} & {春} & {夏} & {秋} & {冬} & {春} & {夏} & {秋} \\
\midrule
\SI{10}{\hecto\pascal}  & 12.8 & 11.7 & 3.0  & 15.5 & 14.0 & 14.7 & 2.5  & 13.1 & 14.2 & 14.8 & 2.5  & 12.9 \\
\SI{30}{\hecto\pascal}  & 10.2 & 11.4 & 4.7  & 11.7 & 11.2 & 12.0 & 2.8  & 9.7  & 11.1 & 12.0 & 2.7  & 9.2 \\
\SI{100}{\hecto\pascal} & 6.3  & 9.9  & 14.7 & 8.2  & 4.4  & 7.4  & 16.4 & 8.1  & 4.2  & 7.2  & 16.0 & 8.2 \\
\SI{200}{\hecto\pascal} & 9.0  & 16.5 & 18.6 & 12.1 & 5.4  & 10.1 & 21.9 & 10.3 & 5.3  & 9.8  & 22.2 & 10.2 \\
\SI{500}{\hecto\pascal} & 20.2 & 26.1 & 19.4 & 18.0 & 9.0  & 13.9 & 16.3 & 12.6 & 8.8  & 13.3 & 16.1 & 12.3 \\
\SI{850}{\hecto\pascal} & 23.2 & 36.1 & 33.3 & 32.6 & 24.2 & 43.4 & 27.3 & 46.3 & 25.3 & 43.8 & 25.3 & 45.4 \\
\bottomrule
\end{tabular}
\end{table}

\SI{10}{\hecto\pascal} と \SI{30}{\hecto\pascal} では夏が 3 領域とも
\ang{2.5} から \ang{4.7} で、他の季節の 3--5 分の 1 である。
夏の成層圏東風はほとんど向きを変えない。

ところが \SI{100}{\hecto\pascal} から \SI{200}{\hecto\pascal} では順序が反転し、
和歌山では夏が \ang{16.4} から \ang{21.9}、冬が \ang{4.4} から \ang{5.4} になる。
愛媛も同じ値に並ぶ。
北海道でも向きは同じだが、冬が \ang{6.3} から \ang{9.0} と和歌山ほど静かではない。

\SI{500}{\hecto\pascal} では北と南の性格がはっきり分かれる。
和歌山と愛媛の冬は \ang{9} 前後で落ち着いているのに、北海道の冬は \ang{20.2} ある。
中層でも北海道は年じゅう向きが振れている。

\begin{whybox}{持続性が高いのに向きが振れるのはなぜか}
この 2 つは矛盾しない。自己相関が測っているのは\textbf{偏差}の持続性で、
風向の前日差が測っているのは\textbf{風向そのもの}の動きである。

夏の \SI{200}{\hecto\pascal} は平均風が弱い。
弱い風に同じ大きさの偏差が乗ると、向きは大きく回る。
冬は平均風が \SI{70}{\metre\per\second} 前後の強いジェットなので、
同じ偏差では向きがほとんど動かない。

落下点を決めるのは向きと強さの両方なので、
運用上は\textbf{風向の表のほうが直接役に立つ}。
\end{whybox}

\SI{850}{\hecto\pascal} ではまた別の順序になり、南の 2 領域では春（\ang{43}--\ang{44}）と秋（\ang{45}--\ang{46}）が大きい。
地表付近は季節の移り変わりそのものが荒いということで、
これも上空とは別の理由による。

\subsubsection{風の分散はどこにあるか}

最後に、風の分散を 3 つに分けた。

\begin{figure}[H]
\centering
\includegraphics[width=0.85\linewidth]{fidelity/G6-wakayama_variance_split.pdf}
\caption[風の分散の内訳]{風の分散を、季節が説明する分・季節内の日々の変動・1 日の中の変動に分けたもの。時期別の気候値で説明できるのは左の帯だけで、右の 2 つはその日の天気で決まる。}
\label{fig:g6-wakayama}
\end{figure}

\begin{figure}[H]
\centering
\includegraphics[width=0.85\linewidth]{fidelity/G6-hokkaido_variance_split.pdf}
\caption[風の分散の内訳]{風の分散を、季節が説明する分・季節内の日々の変動・1 日の中の変動に分けたもの。時期別の気候値で説明できるのは左の帯だけで、右の 2 つはその日の天気で決まる。}
\label{fig:g6-hokkaido}
\end{figure}

\begin{figure}[H]
\centering
\includegraphics[width=0.85\linewidth]{fidelity/G6-ehime_variance_split.pdf}
\caption[風の分散の内訳]{風の分散を、季節が説明する分・季節内の日々の変動・1 日の中の変動に分けたもの。時期別の気候値で説明できるのは左の帯だけで、右の 2 つはその日の天気で決まる。}
\label{fig:g6-ehime}
\end{figure}


\begin{table}[H]
\centering
\caption{風の分散の内訳（\si{\metre\squared\per\second\squared}、括弧内は全体に占める割合）。}
\label{tab:variance-split}
\begin{tabular}{@{}lrrr@{}}
\toprule
 & 北海道 & 和歌山 & 愛媛 \\
\midrule
全体 & 115.1 & 169.0 & 169.7 \\
季節が説明する分 & 20.9（18 \%） & 76.7（45 \%） & 79.2（47 \%） \\
季節内の日々の天気 & 81.5（71 \%） & 81.7（48 \%） & 80.4（47 \%） \\
1 日の中の変化 & 12.7（11 \%） & 10.6（6 \%） & 10.2（6 \%） \\
\bottomrule
\end{tabular}
\end{table}

和歌山と愛媛では、\textbf{気候値で説明できるのは風の変動のおよそ半分まで}で、
残りの半分はその日の天気が決めている。
北海道では季節の分が 18 \% しかなく、7 割を日々の天気が決めている。
これは \ref{sec:validation} 章の結論と同じことを別の角度から言っている。
半月別の表は計画を立てるための道具であって、予報ではない。

\subsection{打ち上げがずれたら落下点はどれだけ動くか}
\label{sec:fidelity-slip}

準備が遅れて N 時間後に上げ直したとき、落下点はどこへ動くか。
これはシミュレータの誤差ではない。
シミュレータはどちらの打ち上げについても正しく答えたはずである。
測っているのは\textbf{予定がずれることのリスク}で、
打ち上げ計画にどれだけ余裕が要るかを決める数字である。

\begin{figure}[H]
\centering
\includegraphics[width=0.92\linewidth]{fidelity/G4-wakayama_slip_curve.pdf}
\caption[打ち上げが遅れたときの落下点の移動]{打ち上げが遅れたとき、予測落下点がどれだけ動くか。実線が中央値、帯が p90 まで。これはシミュレータの誤差ではなく、予定がずれることのリスクである。}
\label{fig:g4-wakayama}
\end{figure}

\begin{figure}[H]
\centering
\includegraphics[width=0.92\linewidth]{fidelity/G4-hokkaido_slip_curve.pdf}
\caption[打ち上げが遅れたときの落下点の移動]{打ち上げが遅れたとき、予測落下点がどれだけ動くか。実線が中央値、帯が p90 まで。これはシミュレータの誤差ではなく、予定がずれることのリスクである。}
\label{fig:g4-hokkaido}
\end{figure}

\begin{figure}[H]
\centering
\includegraphics[width=0.92\linewidth]{fidelity/G4-ehime_slip_curve.pdf}
\caption[打ち上げが遅れたときの落下点の移動]{打ち上げが遅れたとき、予測落下点がどれだけ動くか。実線が中央値、帯が p90 まで。これはシミュレータの誤差ではなく、予定がずれることのリスクである。}
\label{fig:g4-ehime}
\end{figure}


\begin{table}[H]
\centering
\caption{打ち上げが 1 日遅れたときの落下点の移動（\si{\kilo\metre}）。
北海道はどの季節でも和歌山・愛媛より大きい。}
\label{tab:slip}
\sisetup{table-format=3.1}
\begin{tabular}{@{}lSSSSSS@{}}
\toprule
& \multicolumn{2}{c}{北海道} & \multicolumn{2}{c}{和歌山} & \multicolumn{2}{c}{愛媛} \\
\cmidrule(lr){2-3}\cmidrule(lr){4-5}\cmidrule(lr){6-7}
季節 & {中央値} & {p90} & {中央値} & {p90} & {中央値} & {p90} \\
\midrule
冬 & 55.5 & 108.9 & 47.3 & 89.6 & 45.4 & 86.2 \\
春 & 63.2 & 122.5 & 49.8 & 96.7 & 48.1 & 93.7 \\
夏 & 43.0 & 88.6  & 29.7 & 65.5 & 29.1 & 63.4 \\
秋 & 55.9 & 109.6 & 36.4 & 74.5 & 35.1 & 71.9 \\
\bottomrule
\end{tabular}
\end{table}

1 日の遅れで、和歌山は春 \SI{49.8}{\kilo\metre}、冬 \SI{47.3}{\kilo\metre}、
夏はその 6 割の \SI{29.7}{\kilo\metre}。
\textbf{春がもっとも不安定で、夏がもっとも安定}という、前節と同じ順序である。
愛媛はどの季節も和歌山より \SIrange{1}{2}{\kilo\metre} 小さいだけで、同じ順序に並ぶ。

北海道も順序は同じだが、\textbf{どの季節でも和歌山より大きい}。
いちばん安定した夏でも \SI{43.0}{\kilo\metre} 動き、
これは和歌山のいちばん悪い春（\SI{49.8}{\kilo\metre}）に近い。
北海道では\textbf{いちばん安定した季節を選んでも、和歌山のいちばん不安定な季節と同じくらい落下点が動く}。
偏差が続く時間の短さ（表 \ref{tab:decorrelation-both}）がそのまま表れている。

\subsection{シミュレータ自身の近似誤差}
\label{sec:fidelity-error}

\subsubsection{格子の中を実際に飛ばす}

(1) と (2) の大きさを測るには、固定しない場合の答えが要る。
そこで、機体を格子の中で実際に飛ばす積分器を書いた。
水平は緯度経度で双線形、時間は 6 時間解析の間を線形に内挿し、
移動しながら、そのとき自分がいる場所・時刻の風を読む。

ここから 4 通りを走らせて比べる。

\begin{table}[H]
\centering
\caption{比べた 4 通り。差が測りたいものである。}
\label{tab:variants}
\begin{tabular}{@{}llll@{}}
\toprule
呼び方 & 風を読む時刻 & 風を読む場所 & 対応するもの \\
\midrule
基準 & 飛行中の時刻 & 機体の位置 & 再解析で言える最良の答え \\
時刻のみ固定 & 打ち上げ時刻 & 機体の位置 & (1) だけの影響 \\
地点のみ固定 & 飛行中の時刻 & 打ち上げ点 & (2) だけの影響 \\
両方固定 & 打ち上げ時刻 & 打ち上げ点 & \textbf{現行のシミュレータ} \\
\bottomrule
\end{tabular}
\end{table}

\begin{whybox}{「格子間を内挿しない」という方針に反しないのか}
本パイプラインは格子間の内挿を避けてきた。
2 つの p90 を内挿しても、得られる数は何の p90 でもないからである。

ここで内挿しているのは統計量ではなく\textbf{風の場そのもの}で、
これは格子が標本化している連続な物理場である。
スペースバルーンは格子点の上だけを飛ぶわけではないので、
内挿しないほうがかえって誤差になる。

区別を保つため、この双線形内挿は \texttt{stats/trajectory.py} の中だけに置き、
統計を扱うモジュールには一切持ち込んでいない。
\end{whybox}

鉛直方向は、既存の \texttt{drift.layer\_seconds} と同じ層モデルをそのまま使っている。
気圧面のあいだを内挿すればもう少し忠実にはなるが、
そうすると\texttt{両方固定}と既存の閉形式との差に鉛直の精度差が混ざり、
測ろうとしている「固定したことによる誤差」が分離できなくなる。

高度プロファイルは領域で 1 本に統一した。
4 通りがまったく同じ $z(t)$ を飛ばないと、
差の一部が飛行時間の違いになってしまうからである。
実測すると、領域内で気圧面の高度がばらつく幅は最大 \SI{203}{\metre}、
上昇する \SI{30}{\kilo\metre} の 0.68\,\% であった。

\subsubsection{積分器が正しいことをどう確かめたか}

新しい積分器が出す数字は、積分器が正しいときだけ意味を持つ。
そこで答えが独立にわかる入力で固定した。

もっとも強い確認は、\textbf{両方固定}が既存の閉形式
\texttt{drift.displacements} を再現することである。
両方固定すれば積分は既存の総和 $\sum_k u_k T_k$ に厳密に一致するはずで、
実際、東西方向だけの風（緯度が動かないので地図因子が一定になる）で
機械精度まで一致した。一様な風ではなく、高度・場所・時刻すべてに依存する
風で一致するので、層時間・層の番号付け・固定の切り替えのすべてが検証される。

他に、無風で変位ゼロ、定常一様風で球面上の解析解（航程線の積分）を
相対誤差 $10^{-7}$ で再現、時間だけ変わる風では時刻固定にだけ誤差が出て
地点固定には出ない（空間だけ変わる風ではその逆）、
刻み幅 \SIrange{30}{120}{\second} で着地点が \SI{1}{\kilo\metre} 以内に一致、
といった確認を行った。

\subsubsection{結果}

和歌山は \num{1051992} 回（10 年 $\times$ 6 時間ごと $\times$ 72 格子点）、
北海道は \num{8328270} 回（同じく 570 格子点）、
愛媛は \num{1577988} 回（同じく 108 格子点）の打ち上げを 4 通りずつ飛ばした。

\begin{whybox}{北海道では 7 \% の軌道がデータの外へ出る}
和歌山ではマスター箱の外へ出たものが 0 件だった。
北海道では \num{579271} 件、\SI{6.96}{\percent} が出る。

原因は箱の形である。マスター箱の東端は \ang{150}E で、
北海道領域の東端は \ang{149}E。東へ向かう余裕は 1 度、
\SI{81}{\kilo\metre} しかない（北へは \ang{1.5}、\SI{166}{\kilo\metre} ある）。
冬に \SI{200}{\kilo\metre} 東へ飛ぶ軌道はそこを越えてしまう。

脱出はほぼ東向きである。2020 年 1 月に出た軌道の 82 \%、
7 月では 100 \% が、領域の東端から 1 度以内で上げたものだった。
年間を通すと打ち上げ経度は \ang{138.0} から \ang{148.9} に及ぶので、
まれには西寄りの地点からも出る。それでも主因は東端の余裕不足である。

出た軌道は集計から外してある。
ただし\textbf{これは無作為な脱落ではない}。
速い風に乗って東へ走った軌道だけが抜けるので、
外した集団は飛行距離も誤差も平均より大きい。
10 年ぶんで飛行距離 \SI{206}{\kilo\metre} 対 \SI{132}{\kilo\metre}、
両方固定の誤差 \SI{4.87}{\kilo\metre} 対 \SI{4.11}{\kilo\metre} である。

したがって以下の北海道の誤差は\textbf{わずかに小さめに出ている}。
脱落率と差から見積もると、両方固定の中央値で \SI{1.3}{\percent}、
地点のみ固定で \SI{2.2}{\percent}、時刻のみ固定で \SI{1.1}{\percent}。
どの結論も動かない大きさである。
実装は \texttt{fidelity\_out.exclusion\_bias} が毎回この量を計算して記録する。
箱を広げれば消える制約だが、それには再取得が要る。
\end{whybox}

\begin{whybox}{愛媛では 6 \% の軌道が西端から出る}
愛媛でも \num{93806} 件、\SI{5.94}{\percent} がマスター箱の外へ出た。
こちらは北海道と逆向きで、\textbf{西端}（\ang{132}E）から出ている。
愛媛の領域は箱の西端に接しているので（第 \ref{sec:ehime} 章）、
\ang{132.0}--\ang{133.5}E から上げて夏に西へ流れた軌道が、すぐに外へ出る。

北海道と違い、外れた集団は飛行距離が\textbf{短い}
（\SI{95}{\kilo\metre} 対 \SI{152}{\kilo\metre}）。
夏の弱い風に乗った飛行だからで、固定による誤差も平均より小さい。
したがって以下の愛媛の誤差は\textbf{わずかに大きめに出ている}。
大きさは両方固定の中央値で \SI{1.1}{\percent}、
地点のみ固定で \SI{1.8}{\percent}、時刻のみ固定で \SI{0.8}{\percent} で、
やはり結論は動かない。
\end{whybox}

\begin{figure}[H]
\centering
\includegraphics[width=1.0\linewidth]{fidelity/G5-wakayama_error_breakdown.pdf}
\caption[風を固定したことによる誤差]{基準軌道の着地点からのずれ。棒が中央値、細線が p90 まで。左は \si{\kilo\metre}、右は飛行距離に対する割合。冬は飛行距離が夏の 4 倍あるので、2 つのパネルは季節を逆に並べる。どちらか一方にまとめると、測定ではなく判断になってしまう。}
\label{fig:g5-wakayama}
\end{figure}

\begin{figure}[H]
\centering
\includegraphics[width=1.0\linewidth]{fidelity/G5-hokkaido_error_breakdown.pdf}
\caption[風を固定したことによる誤差]{基準軌道の着地点からのずれ。棒が中央値、細線が p90 まで。左は \si{\kilo\metre}、右は飛行距離に対する割合。冬は飛行距離が夏の 4 倍あるので、2 つのパネルは季節を逆に並べる。どちらか一方にまとめると、測定ではなく判断になってしまう。}
\label{fig:g5-hokkaido}
\end{figure}

\begin{figure}[H]
\centering
\includegraphics[width=1.0\linewidth]{fidelity/G5-ehime_error_breakdown.pdf}
\caption[風を固定したことによる誤差]{基準軌道の着地点からのずれ。棒が中央値、細線が p90 まで。左は \si{\kilo\metre}、右は飛行距離に対する割合。冬は飛行距離が夏の 4 倍あるので、2 つのパネルは季節を逆に並べる。どちらか一方にまとめると、測定ではなく判断になってしまう。}
\label{fig:g5-ehime}
\end{figure}


\begin{table}[H]
\centering
\caption{基準軌道からのずれ（中央値、\si{\kilo\metre}）。
「現行」が実際に動いているシミュレータにあたる。}
\label{tab:fidelity-error}
\sisetup{table-format=1.1}
\begin{tabular}{@{}lSSSS[table-format=3.0]S@{}}
\toprule
季節 & {時刻のみ} & {地点のみ} & {両方（現行）} & {飛行距離 [\si{\kilo\metre}]} & {現行の割合 [\%]} \\
\midrule
\multicolumn{6}{@{}l}{\textbf{北海道}} \\
冬 & 3.6 & 5.8 & 5.1 & 176 & 3.1 \\
春 & 4.1 & 4.7 & 4.2 & 134 & 3.4 \\
夏 & 3.0 & 2.7 & 3.2 & 73  & 4.6 \\
秋 & 3.7 & 5.3 & 4.5 & 164 & 2.9 \\
\addlinespace
\multicolumn{6}{@{}l}{\textbf{和歌山}} \\
冬 & 3.3 & 6.1 & 5.3 & 267 & 2.0 \\
春 & 3.5 & 4.1 & 3.9 & 165 & 2.5 \\
夏 & 2.2 & 2.0 & 2.8 & 63  & 4.6 \\
秋 & 2.7 & 3.2 & 3.5 & 125 & 3.0 \\
\addlinespace
\multicolumn{6}{@{}l}{\textbf{愛媛}} \\
冬 & 3.2 & 5.9 & 5.2 & 267 & 2.0 \\
春 & 3.4 & 4.0 & 3.9 & 165 & 2.5 \\
夏 & 2.2 & 1.9 & 2.8 & 63  & 4.6 \\
秋 & 2.6 & 3.0 & 3.4 & 124 & 3.0 \\
\bottomrule
\end{tabular}
\end{table}

まず絶対値でみると、誤差は 3 領域・4 季節のどれでも
\SIrange{2.8}{5.3}{\kilo\metre} に収まり、飛行距離の \SIrange{2}{5}{\percent} である。
前節の「1 日ずれたら \SIrange{29}{63}{\kilo\metre} 動く」に比べると 1 桁小さい。

\begin{whybox}{風を固定する近似は、直しても精度はほとんど上がらない}
現行シミュレータの近似誤差は、打ち上げが半日ずれたときの移動の
1 割前後にすぎない。
測れなかった予報誤差 (4) はさらに大きいと考えられる。

風を 4 次元で引き直す改修は \texttt{wind.py} だけの話なので難しくはないが、
得られるのは中央値で数キロメートルである。
落下点の精度を上げたいなら、
もっと大きい誤差源から直すほうが効果が大きい。

\textbf{この結論は 3 領域で共通である。} 気候がまるで違う北海道と南の 2 領域で同じ
\SIrange{2}{5}{\percent} に収まったことは、この数字が土地の事情ではなく
飛行時間（約 \SI{137}{\minute}）と風の時間・空間スケールの比で決まっていることを示唆する。
\end{whybox}

\subsubsection{冬には 2 つの固定が打ち消し合う}

表をよく見ると、和歌山の冬だけ\textbf{両方固定（\SI{5.3}{\kilo\metre}）が
地点のみ固定（\SI{6.1}{\kilo\metre}）より小さい}。
誤差を 1 つ足したのに、誤差が減っている。

愛媛の冬も \SI{5.2}{\kilo\metre} 対 \SI{5.9}{\kilo\metre} で同じである。
北海道でも同じことが起きている。冬は \SI{5.1}{\kilo\metre} 対 \SI{5.8}{\kilo\metre}、
秋は \SI{4.5}{\kilo\metre} 対 \SI{5.3}{\kilo\metre} で、
どちらも両方固定のほうが小さい。

これは偶然ではなく、冬の流れの性質から説明がつく。
冬は強い西風で、気圧の谷や尾根も西から東へ動いていく。
機体も東へ運ばれる。

地点を固定すると、機体が着くはずの東側の風ではなく、打ち上げ点の風を使う。
時刻を固定すると、到着時刻の風ではなく、打ち上げ時刻の風を使う。
パターン自体が東へ移動しているので、
「東側の、あとの時刻の風」は「打ち上げ点の、いまの風」に似ている。
2 つのずれが互いを部分的に埋める。

和歌山の夏はこうならない。
両方固定 \SI{2.8}{\kilo\metre} が時刻のみ \SI{2.2}{\kilo\metre} と
地点のみ \SI{2.0}{\kilo\metre} のどちらより大きく、2 つの誤差がそのまま足し合わさる。
夏の和歌山には東進する強い流れがないので、埋め合わせが起きない。
愛媛の夏も \SI{2.8}{\kilo\metre} 対 \SI{2.2}{\kilo\metre}・\SI{1.9}{\kilo\metre} で、同じく足し合わさる。

\textbf{北海道の夏は、そのどちらでもない。}

季節をまとめたことで生じた見かけではないことを確かめるため、
北海道と和歌山について、月ごとに分けて計算し直した。

\begin{table}[H]
\centering
\caption{両方固定の誤差を、大きいほうの片側固定で割った比。
1 を下回れば打ち消し合い、上回れば足し合わせである。}
\label{tab:cancellation}
\sisetup{table-format=1.2}
\begin{tabular}{@{}lSSSS@{}}
\toprule
月 & {時刻のみ [\si{\kilo\metre}]} & {地点のみ [\si{\kilo\metre}]}
   & {両方 [\si{\kilo\metre}]} & {比} \\
\midrule
\multicolumn{5}{@{}l}{\textbf{北海道}} \\
2018年12月 & 3.96 & 7.10 & 5.92 & 0.83 \\
2019年1月  & 3.98 & 6.55 & 5.85 & 0.89 \\
2019年2月  & 4.20 & 6.83 & 6.07 & 0.89 \\
\addlinespace
2019年6月  & 3.23 & 2.63 & 2.95 & 0.91 \\
2019年7月  & 2.58 & 1.69 & 2.53 & 0.98 \\
2019年8月  & 2.66 & 3.12 & 3.09 & 0.99 \\
\addlinespace
\multicolumn{5}{@{}l}{\textbf{和歌山}} \\
2018年12月 & 3.22 & 6.25 & 5.64 & 0.90 \\
2019年1月  & 3.53 & 8.27 & 6.68 & 0.81 \\
2019年2月  & 3.54 & 6.17 & 5.59 & 0.91 \\
\addlinespace
2019年6月  & 3.33 & 3.49 & 4.06 & 1.16 \\
2019年7月  & 2.21 & 1.80 & 2.56 & 1.16 \\
2019年8月  & 2.11 & 2.08 & 2.74 & 1.30 \\
\bottomrule
\end{tabular}
\end{table}

\textbf{冬は 2 領域とも打ち消し合う。} 6 か月すべてが 1 を下回り、
範囲も \numrange{0.81}{0.91} とよく揃っている。
季節をまとめたことによる見かけではない。

\textbf{夏は 2 領域で違う。}
和歌山の夏は 3 か月とも 1 を上回り（\numrange{1.16}{1.30}）、単純な足し合わせである。
北海道の夏は \numrange{0.91}{0.99} で、1 をわずかに下回るところに貼りついている。
足し合わせでも打ち消し合いでもない。

\begin{whybox}{北海道の夏に「中間」が出るのはなぜか}
打ち消し合いは、擾乱が機体と同じ向き・同じくらいの速さで動くときに起きる。

和歌山の夏は亜熱帯ジェットが北へ抜けていて、動く擾乱がほとんどない。
だから 2 つのずれは独立に効き、足し合わさる。

北海道の夏は違う。ジェットは弱まるが消えはせず、
\SI{200}{\hecto\pascal} で平均 \SI{29.9}{\metre\per\second} 残り、
低気圧も通り続ける（\SI{200}{\hecto\pascal} の偏差が \num{1.8} 日しか続かないのはそのためである）。
移動する擾乱が弱いながらも存在するので、打ち消しが部分的にしか働かない。

つまりこの比は、\textbf{その季節・その場所に「東へ動くパターン」が
どれだけあるか}を測る指標として読める。
冬の 0.8 台、和歌山の夏の 1.2 台、北海道の夏の 0.95 前後という並びは、
そのまま擾乱の活発さの順序になっている。
\end{whybox}

\begin{figure}[H]
\centering
\includegraphics[width=1.0\linewidth]{fidelity/G7-01-wakayama_trajectory_map.pdf}
\caption[冬の 1 例]{2020 年 1 月 15 日 00 UTC に和歌山の中心格子から上げた場合の飛行経路。左が全体、右が着地点の拡大。地点のみ固定（橙の破線）と両方固定（黄）が南側にそろい、基準（黒）と時刻のみ固定（青の破線）が北側にそろう。冬は地点の固定による誤差のほうが大きいことが、1 例からも見てとれる。}
\label{fig:g7-01-wakayama}
\end{figure}

\begin{figure}[H]
\centering
\includegraphics[width=1.0\linewidth]{fidelity/G7-07-wakayama_trajectory_map.pdf}
\caption[夏の 1 例]{同じ図の 7 月版。飛行距離が \SI{80}{\kilo\metre} ほどと短く、4 通りの差も \SI{2.5}{\kilo\metre} にとどまる。拡大パネルの縮尺が冬（\SI{14}{\kilo\metre}）と夏（\SI{2}{\kilo\metre}）で違う点に注意。}
\label{fig:g7-07-wakayama}
\end{figure}

\begin{figure}[H]
\centering
\includegraphics[width=1.0\linewidth]{fidelity/G7-01-hokkaido_trajectory_map.pdf}
\caption[冬の 1 例（北海道）]{2020 年 1 月 15 日 00 UTC に VOISING から上げた場合の飛行経路。
左が全体、右が着地点の拡大。
この例では 4 通りの差が \SI{2.9}{\kilo\metre} しかなく、
和歌山の同じ日（図\ref{fig:g7-01-wakayama}、\SI{13.1}{\kilo\metre}）よりずっと小さい。
和歌山では基準（黒）と時刻のみ固定（青）、
地点のみ固定（橙）と両方固定（黄）がそれぞれ対になるのに対し、
ここでは 4 本が順に並ぶだけで、そうした対応は見えない。
\textbf{1 例は 1 例であって統計ではない}。
どちらが典型かは表\ref{tab:fidelity-error} で判断してほしい。}
\label{fig:g7-01-hokkaido}
\end{figure}

\begin{figure}[H]
\centering
\includegraphics[width=1.0\linewidth]{fidelity/G7-07-hokkaido_trajectory_map.pdf}
\caption[夏の 1 例（北海道）]{同じ図の 7 月版。
拡大パネルの縮尺が冬と夏で違う点に注意。}
\label{fig:g7-07-hokkaido}
\end{figure}
\begin{figure}[H]
\centering
\includegraphics[width=1.0\linewidth]{fidelity/G7-01-ehime_trajectory_map.pdf}
\caption[冬の 1 例]{2020 年 1 月 15 日 00 UTC に南レク御荘公園から上げた場合の飛行経路。左が全体、右が着地点の拡大。和歌山の同じ日（図\ref{fig:g7-01-wakayama}）と同じく、地点のみ固定（橙の破線）と両方固定（黄）が南側に、基準（黒）と時刻のみ固定（青の破線）が北側にそろう。現行との差は \SI{11.1}{\kilo\metre}。}
\label{fig:g7-01-ehime}
\end{figure}

\begin{figure}[H]
\centering
\includegraphics[width=1.0\linewidth]{fidelity/G7-07-ehime_trajectory_map.pdf}
\caption[夏の 1 例]{同じ図の 7 月版。いったん東へ出てから西へ戻り、また東へ流れる経路で、4 通りの差は \SI{2.7}{\kilo\metre} にとどまる。拡大パネルの縮尺が冬（\SI{8}{\kilo\metre}）と夏（\SI{4}{\kilo\metre}）で違う点に注意。拡大パネルの左下に見える線は、飛行の途中で着地点の近くを一度通った部分である。}
\label{fig:g7-07-ehime}
\end{figure}


\subsection{結論：どの季節がいちばんシミュレーション通りか}
\label{sec:fidelity-answer}

答えは、何を「通り」と呼ぶかで変わる。そこで両方を出す。
そして\textbf{北海道と南の 2 領域で答えが違う}。和歌山と愛媛の答えは同じである。

\begin{table}[H]
\centering
\caption{総合。絶対値と割合で順位が逆になり、その逆転の仕方が北と南で違う。}
\label{tab:fidelity-ranking}
\sisetup{table-format=2.1}
\begin{tabular}{@{}lS[table-format=3.0]SS[table-format=1.1]SS[table-format=3.0]@{}}
\toprule
季節 & {飛行距離 [\si{\kilo\metre}]} & {近似誤差 [\si{\kilo\metre}]}
     & {割合 [\%]} & {1 日ずれ [\si{\kilo\metre}]} & {同 p90} \\
\midrule
\multicolumn{6}{@{}l}{\textbf{北海道}} \\
冬 & 176 & 5.1 & 3.1 & 55.5 & 109 \\
春 & 134 & 4.2 & 3.4 & 63.2 & 123 \\
夏 & 73  & 3.2 & 4.6 & 43.0 & 89 \\
秋 & 164 & 4.5 & 2.9 & 55.9 & 110 \\
\addlinespace
\multicolumn{6}{@{}l}{\textbf{和歌山}} \\
冬 & 267 & 5.3 & 2.0 & 47.3 & 90 \\
春 & 165 & 3.9 & 2.5 & 49.8 & 97 \\
夏 & 63  & 2.8 & 4.6 & 29.7 & 66 \\
秋 & 125 & 3.5 & 3.0 & 36.4 & 75 \\
\addlinespace
\multicolumn{6}{@{}l}{\textbf{愛媛}} \\
冬 & 267 & 5.2 & 2.0 & 45.4 & 86 \\
春 & 165 & 3.9 & 2.5 & 48.1 & 94 \\
夏 & 63  & 2.8 & 4.6 & 29.1 & 63 \\
秋 & 124 & 3.4 & 3.0 & 35.1 & 72 \\
\bottomrule
\end{tabular}
\end{table}

\paragraph{絶対的な精度がほしいなら、3 領域とも夏。}
近似誤差が最小（和歌山と愛媛 \SI{2.8}{\kilo\metre}、北海道 \SI{3.2}{\kilo\metre}）で、
飛行距離も最短、1 日ずれたときの移動も最小である。
和歌山と愛媛の夏は \SI{200}{\hecto\pascal} の偏差が \numrange{10.7}{10.8} 日続くので、
前日の観測が翌日も生きる。
北海道の夏は偏差が \num{1.8} 日しか続かないが、それでも 4 季節の中では最良である。

\paragraph{比例的に忠実な季節は、領域で分かれる。}
和歌山と愛媛は\textbf{冬}（どちらも \SI{2.0}{\percent}）。
強いジェットで向きの振れも \SI{200}{\hecto\pascal} で \ang{5.4} と小さく、
「どちらへ飛ぶか」がもっとも予測しやすい。

北海道は\textbf{秋}（\SI{2.9}{\percent}）で、冬（\SI{3.1}{\percent}）はわずかに劣る。
差は小さいが、和歌山のように冬が突出して良いということはない。
北海道の冬のジェットは和歌山ほど安定していないためで、
\SI{500}{\hecto\pascal} の風向が冬でも前日から \ang{20.2} 振れる
（和歌山は \ang{9.0}）。

\paragraph{春は 3 領域とも避けたほうがよい。}
1 日ずれたときの移動が和歌山 \SI{49.8}{\kilo\metre}、愛媛 \SI{48.1}{\kilo\metre}、北海道 \SI{63.2}{\kilo\metre} で、
どちらも 4 季節中最大である。
\SI{200}{\hecto\pascal} の偏差が続く時間も最短（和歌山 \num{1.7} 日、愛媛 \num{1.9} 日、北海道 \num{1.2} 日）。
\SI{850}{\hecto\pascal} の風向は前日から \ang{36} から \ang{43} 振れる。

\paragraph{そして領域どうしを比べると、北海道のほうが一貫して扱いにくい。}
1 日ずれたときの移動は\textbf{どの季節でも和歌山・愛媛より大きく}、
北海道のいちばん良い夏（\SI{43.0}{\kilo\metre}）が
和歌山のいちばん悪い春（\SI{49.8}{\kilo\metre}）に近い。
飛行距離は和歌山の冬のほうが長い（\SI{267}{\kilo\metre} 対 \SI{176}{\kilo\metre}）のに、
それでも移動は和歌山のほうが小さい。
距離ではなく\textbf{風の入れ替わりの速さ}が差を生んでいる。

\paragraph{まとめると}、狭い範囲で確実に回収したいなら 3 領域とも夏。
方向を読みたいなら和歌山と愛媛は冬、北海道は秋。春は避ける。
そして北海道では、季節を選んでも和歌山・愛媛ほどの余裕は得られない。

なお繰り返しになるが、これらはすべて
「予報が完璧だったとして」の話である（\ref{sec:fidelity-scope} 節）。
実際の運用では予報誤差 (4) が上乗せされ、本章の数字より外れる。

\subsection{副産物：既存の閉形式にある地図因子のずれ}
\label{sec:fidelity-mapfactor}

積分器と閉形式を突き合わせる過程で、既存の \texttt{drift.displacements} に
系統的なずれがあることがわかった。報告しておく。

閉形式は $\Delta x = \sum_k u_k T_k$ を直接メートルで出していて、
式のどこにも緯度が現れない。
一方で実際の機体は飛行中に南北にも動き、
経度 1 度あたりの距離を決める $\cos(\text{緯度})$ が飛行中に変わる。

このずれは、南北移動が大きい冬の飛行でいちばん大きくなる。
2020 年 1 月で測ると中央値 \SI{0.44}{\kilo\metre}、p90 で \SI{1.12}{\kilo\metre}、
最大 \SI{1.93}{\kilo\metre}。10 年ぶんの最大は \SI{3.45}{\kilo\metre} であった。
7 月は飛行距離が短いので中央値 \SI{0.07}{\kilo\metre} にとどまる。

第 \ref{sec:hokkaido} 章から第 \ref{sec:ehime} 章の落下位置は
すべてこの閉形式で計算しているので、
そこに載せた平均落下位置には冬でおよそ \SI{0.4}{\kilo\metre}、
最悪で数キロメートルの系統誤差が含まれている。
気候値としての比較（どの時期がどちらへ飛ぶか）には影響しない大きさだが、
1 回の飛行の予測に使うには無視できない。
本章の積分器は球面の幾何を正しく扱っているので、この誤差を持たない。
